%Abstract%
\section*{\centering{ABSTRACT}} 
\noindent Technical Debt (TD) in documentation is a challenge in the software industry. In the Microentrepreneur Support System (SFM), an empirical development flow resulted in TD accumulation, causing communication failures and hindering tests. This work aims to report the analysis of this TD and propose its mitigation through formal artifacts. The methodology encompassed requirements elicitation, the construction of a Traceability Matrix, and a questionnaire administered to seven developers to assess the system's perceived completeness. Results revealed that 82\% of the mapped flows lacked test coverage. Furthermore, the survey refuted the research's initial empirical observation: a module previously assessed as consolidated revealed, from the team's perspective, partial implementation and high comprehension divergence. Additionally, it was found that the developers acknowledge the high potential of the produced documentation to assist them in their daily activities. It is concluded that formal documentation provides a solid structural foundation, mitigating the diagnosed opacity, guiding coding, and enabling future test automation, thus attesting to the importance of Requirements Engineering for software quality.

\noindent{\bf Key-words:} $<$Requirements Engineering$>$, $<$Software Engineering$>$, $<$Software Testing$>$, $<$Technical Debt$>$, $<$Traceability Matrix$>$.